\documentclass[UTF8,11pt]{ctexart}
\usepackage[a4paper,margin=22mm]{geometry}
\usepackage{amsmath,amssymb,bm,graphicx,adjustbox,fvextra,longtable,array}
\xeCJKDeclareCharClass{CJK}{"2160 -> "217F, "2460 -> "249B, "2200 -> "22FF}
\newcommand{\fitmath}[1]{\begin{adjustbox}{max width=\linewidth}$\displaystyle #1$\end{adjustbox}}
\setlength{\parindent}{0pt}\setlength{\parskip}{.65em}\renewcommand{\arraystretch}{1.3}
\sloppy\allowdisplaybreaks
\begin{document}
\section*{2025 年计算机统考 408 · 答案与解析}
\subsection*{2025年全国硕士研究生招生考试计算机学科专业基础试题参考答案}

\subsection*{一、单项选择题}

\begin{center}\begin{adjustbox}{max width=\linewidth}\begin{tabular}{|>{\raggedright\arraybackslash}p{\dimexpr(\linewidth-10\tabcolsep-6\arrayrulewidth)/5\relax}|>{\raggedright\arraybackslash}p{\dimexpr(\linewidth-10\tabcolsep-6\arrayrulewidth)/5\relax}|>{\raggedright\arraybackslash}p{\dimexpr(\linewidth-10\tabcolsep-6\arrayrulewidth)/5\relax}|>{\raggedright\arraybackslash}p{\dimexpr(\linewidth-10\tabcolsep-6\arrayrulewidth)/5\relax}|>{\raggedright\arraybackslash}p{\dimexpr(\linewidth-10\tabcolsep-6\arrayrulewidth)/5\relax}|}\hline
1. B & 2. D & 3. D & 4. B & 5. D\\\hline
6. C & 7. C & 8. C & 9. A & 10. D\\\hline
11. A & 12. D & 13. D & 14. D & 15. C\\\hline
16. B & 17. C & 18. B & 19. A & 20. A\\\hline
21. B & 22. A & 23. D & 24. C & 25. C\\\hline
26. B & 27. D & 28. C & 29. C & 30. A\\\hline
31. C & 32. B & 33. B & 34. C & 35. C\\\hline
36. C & 37. B & 38. A & 39. B & 40. A\\\hline\end{tabular}\end{adjustbox}\end{center}

\subsection*{二、综合应用题}

41．【答案要点】（1）算法的基本设计思想：从后向前扫描一趟数组A，对每个\(\displaystyle A[i]\)（\(\displaystyle 0\le i\le n-1\)），分别找到从\(\displaystyle A[n-1]\)到\(\displaystyle A[i]\)中的最大值Max和最小值Min，然后分以下情况处理：\textcircled{1}若\(\displaystyle A[i]\ge0\)，则\(\displaystyle A[i]\)与Max相乘；\textcircled{2}若\(\displaystyle A[i]<0\)，则\(\displaystyle A[i]\)与Min相乘。相乘的结果保存在\(\displaystyle res[i]\)中。


41．（续）（2）算法实现：


\begin{Verbatim}[breaklines,breakanywhere,fontsize=\small]
void calMulMax(int A[], int res[], int n)
{
    int i, Max, Min;
    Max=Min=A[n-1];
    for(i=n-1; i>=0; i--) {  // 从后向前扫描数组 A
        if(A[i]>Max) Max=A[i];
        else if(A[i]<Min) Min=A[i];
        if(A[i]>=0) res[i]=A[i]*Max;  // 根据 A[i] 的正负性分别处理
        else res[i]=A[i]*Min;
    }
}
\end{Verbatim}

（3）算法的时间复杂度为\(\displaystyle O(n)\)，空间复杂度为\(\displaystyle O(1)\)。


42．【答案要点】（1）完成该工程的最短时间是12。关键活动是a、e、m、n。（2）与e同时进行的活动可能有b、c、d。（3）时间余量最大的活动是j。时间余量是6。（4）b的持续时间最多是4。压缩k的持续时间。


43．【答案要点】（1）Cache组号字段占6位；块内地址字段占6位。虚拟地址中\(\displaystyle \mathrm{VA}_{11}\sim\mathrm{VA}_{6}\)可作为Cache索引。（2）\(\displaystyle d[100]\)的虚拟地址为0180 01B0H。Cache组号是\(\displaystyle 000110\mathrm{B}=6\)。


43．（续）（3）\(\displaystyle d[0]\)在其所在主存块内的偏移量为20H。共发生129次Cache缺失，访问总次数为\(\displaystyle 2048\times2=4096\)，Cache缺失率为\(\displaystyle 129/4096\approx3.15\%\)；数组元素的平均访问时间约为\(\displaystyle 2+3.15\%\times200=8.3\)个时钟周期。（4）数组d分布在3页中。访问数组d所引起的缺页次数是3。


44．【答案要点】（1）R中初始内容是FFFF FFFFH；Q中初始内容是8765 4321H；Y中初始内容是0000 00FFH。控制逻辑部件中包含计数器。由ALUop所控制的ALU运算有加运算、减运算。（2）当x为0000 0000H时，发生除0异常；当d[i]为8000 0000H且x为FFFF FFFFH时，发生除运算溢出异常。CPU检测到除法异常后，需要将断点和程序状态保存到内核栈或者特定寄存器中，关中断，最后跳转到内核中的除法异常处理程序执行。


45．【答案要点】设置信号量及其初值：


\begin{Verbatim}[breaklines,breakanywhere,fontsize=\small]
Semaphore position=3;  // 可挖树坑的数量
Semaphore pit=0;       // 树坑的数量
Semaphore tree=0;      // 待浇水的树苗数
Semaphore mutex=1;     // 实现铁锹的互斥使用
\end{Verbatim}

甲：


\begin{Verbatim}[breaklines,breakanywhere,fontsize=\small]
for(;;) {
    wait(position);
    wait(mutex);
    挖树坑;
    signal(mutex);
    signal(pit);
    ...
}
\end{Verbatim}

乙：


\begin{Verbatim}[breaklines,breakanywhere,fontsize=\small]
for(;;) {
    wait(pit);
    放树苗;
    signal(position);
    wait(mutex);
    填土;
    signal(mutex);
    signal(tree);
    ...
}
\end{Verbatim}

丙：


\begin{Verbatim}[breaklines,breakanywhere,fontsize=\small]
for(;;) {
    wait(tree);
    浇水;
    ...
}
\end{Verbatim}

46．【答案要点】（1）进程控制块位于内核区。该进程处于阻塞态。（2）main()函数的代码位于只读代码段。scanf()和printf()的功能需要通过执行驱动程序实现。（3）ptr被分配在可读写数据段中。length会被分配在用户栈中。ptr指向的字符串位于运行时堆中。


47．【答案要点】（1）单向传播时延为


\[\fitmath{\displaystyle 2\times\dfrac{\displaystyle 36\,000\times10^3}{\displaystyle 300\,000\times10^3}=240\,\mathrm{ms}}\]

最大吞吐量为200 kb/s。最少时间为


\[\fitmath{\displaystyle 240\times10^{-3}+\dfrac{\displaystyle 4000\times8}{\displaystyle 200\times10^3}=400\,\mathrm{ms}}\]

（2）令\(\displaystyle W_s\)为发送窗口，\(\displaystyle k\)为序号字段位数，则有


\[\fitmath{\displaystyle \dfrac{\displaystyle W_s\times\dfrac{\displaystyle 1500\times8}{\displaystyle 200\times10^3}}{\displaystyle 2\times240\times10^{-3}+\dfrac{\displaystyle 1500\times8}{\displaystyle 200\times10^3}}\ge0.8\qquad\text{①}}\]

解得\(\displaystyle W_s\ge7.2\)，向上取整得\(\displaystyle W_s\ge8\)，即发送窗口至少为8。对于GBN，应满足\(\displaystyle 2^k\ge W_s+1\)，即\(\displaystyle 2^k\ge8+1=9\)，\(\displaystyle k\)向上取整得\(\displaystyle k\ge4\)，即序号字段至少需要4位。


（3）作业区子网的子网地址为10.10.10.64/26；管理区子网的子网地址为10.10.10.0/26；生活区子网的子网地址为10.10.10.128/25。

\end{document}
